\begin{table}[htbp]
\centering
\caption{单阶段稠密检索在工程认证文档场景下的典型局限性示例}
\label{tab:dense_limit_cases}
\small
\setlength{\tabcolsep}{4pt}
\renewcommand{\arraystretch}{1.25}
\begin{tabular}{p{0.14\linewidth} p{0.19\linewidth} p{0.26\linewidth} p{0.20\linewidth} p{0.14\linewidth}}
\hline
局限性类型 & 查询内容 & 单阶段稠密检索表现 & 证据缺陷分析 & 双阶段改进效果 \\
\hline
语义相关但证据不足 &
试卷命题与前两届试题的重复率限制是多少？ &
Top-1 返回附件1-9第3个切块（第1页），内容为“ A、B 两套试卷中重复题目不得超过30\% ”；真正对应“前两届试题”的附件1-9第2个切块仅排第2。 &
单阶段检索命中了同一制度文本内的相邻条款，虽与“重复率 30\%”语义高度相关，但约束对象并非“前两届试题”，因此首位证据不能直接支撑回答。 &
重排序后，附件1-9第2个切块（第1页）升至 Top-1，首条证据即可直接支撑“不得超过30\%”这一结论。 \\
\hline
跨块/跨页证据割裂 &
毕业要求和课程体系的合理性评价周期是几年？ &
Top-1 返回附件1-14第3个切块（第1页），给出“学生学满四年毕业时完成达成度评价”；Top-2 返回附件1-5第6个切块（第2页），给出“全校性修订四年一次、课程体系每年可局部调整”；真正关键的附件1-15第7个切块在单阶段仅排第3。 &
问题的完整答案需要同时结合“毕业要求合理性评价”与“课程体系局部调整”两部分证据。单阶段虽有召回，但关键证据分散在不同 chunk 与不同页，前列结果被相近但不充分的“四年周期”条款占据。 &
重排序后，附件1-15第7个切块（第2页）升至 Top-1，附件1-15第14个切块（第3页）进入前列，证据链由分散转为集中，更利于直接作答。 \\
\hline
高召回下噪声增多 &
开卷考试中直接从教材找答案的题目分值限制是多少？ &
Top-1 返回附件1-7第11个切块（第6页），内容为“教材编写绩效最高不超过600分”；Top-2 返回课程规模计分规则；正确证据附件1-9第2个切块（第1页）仅排第4。 &
“教材”“分值”“限制”等表层词项使大量绩效计分条款进入前列，导致真正与开卷命题规则相关的证据被噪声淹没。 &
重排序后，附件1-9第2个切块（第1页）升至 Top-1，直接命中“可从教材找到答案的题目分值不应超过总分值的20\%”。 \\
\hline
\end{tabular}
\end{table}
